% This template was provided for the 2019 Recording Project class at UWL.
% It is designed to provide a simple, low-configuration starter template
% for academic writing with LaTeX and Markdown.
% 
% This template is set up to use the agsm Harvard referencing style by default.

\documentclass[12pt,a4paper, oneside]{report}
\usepackage[top=2cm, bottom=2cm, left=2cm, right=2cm]{geometry}
\usepackage[export]{adjustbox}

%\usepackage[german]{babel}
\usepackage[english]{babel}
\usepackage[utf8]{inputenc}
\usepackage{amsmath}
\usepackage{fancyhdr, graphicx}
\usepackage[T1]{fontenc}
\usepackage[utf8]{inputenc}
\usepackage{lmodern}
\usepackage{csquotes}
\usepackage{natbib}
\usepackage[footnotes,definitionLists,hashEnumerators,smartEllipses,hybrid,citations]{markdown}
\usepackage{amsmath}
\usepackage{longtable}
\usepackage{listings}
\usepackage{color}

\usepackage{float}
\usepackage{wasysym} % für die benutzten Symbole
\usepackage[most]{tcolorbox}


\usepackage{relsize}

\usepackage[most]{tcolorbox}
\newtcolorbox{mybox}[2][]{%
  attach boxed title to top left
               = {yshift=-8pt},
  colback      = blue!5!white,
  colframe     = blue!75!black,
  fonttitle    = \bfseries,
  colbacktitle = blue!85!black,
  title        = #2,#1,
  enhanced,
}

\newtcolorbox{gebox}[2][]{%
  attach boxed title to top left
               = {yshift=-8pt},
  colback      = yellow!5!white,
  colframe     = yellow!75!black,
  fonttitle    = \bfseries,
  colbacktitle = yellow!85!black,
  title        = #2,#1,
  enhanced,
}


\newtcolorbox{nabox}[2][]{%
  attach boxed title to top left
               = {yshift=-8pt},
  colback      = red!5!white,
  colframe     = red!75!black,
  fonttitle    = \bfseries,
  colbacktitle = red!85!black,
  title        = #2,#1,
  enhanced,
}



\setcounter{tocdepth}{4}
\setcounter{secnumdepth}{4}

\setlength{\parskip}{0.5em}

\newcommand{\hash}{\#} % Hashes are difficult in the Markdown environment. typing \hash within a Markdown document will give you a working hash instead.

% A basic Harvard style is the agsm style. Other similar styles are apalike and lsalike.


\bibliographystyle{ieeetr}
%\bibliographystyle{agsm}
\setkeys{Gin}{width=.95\linewidth}
\pagenumbering{gobble}% Remove page numbers (and reset to 1)

% Add your report title here:


\title{
%\vspace{-3cm}
%\hspace{10cm}
%\includegraphics[width=0.33\textwidth]{img/ikarus_logo.png}
Industrial Cyber Security}


% Add your name here:
\author{Herbert Dirnberger}

% The command below auto-generates today's date.
\date{\today}


% The \begin{} command tells LaTeX that we're starting the document.

\begin{document}

% \maketitle prints the title (as defined in the previous section) to the page.

\maketitle

%\textit{Note: The information presented in this document is derived from notes and educational materials of the author, supplemented with information from the provided links. The summary has been generated with the assistance of ChatGPT for the purpose of condensing and simplifying the content.}

% \newpage is fairly self-explanatory!
% \newpage

% \tableofcontents reads the whole document and auto-generates a table of contents based on your section headings.

\setlength{\parindent}{0pt} 

%\paragraph{Executive Summary}
%By applying this 

%\paragraph{Target Group and Scope}
%This document is primarily intended for asset owners, system integrators and service operators in all sections of the asset life cycle: Requirement Engineering, Project planning, Procurement, System Integration and System Service and Operation

\tableofcontents

%\subsection{% \listoffigures generates a list of images used in the document.}
%\listoffigures

\pagenumbering{arabic}% Arabic page numbers (and reset to 1)

% This is where your work goes!
% The \markdownInput{} command inserts the contents of a Markdown file into the document body. The first file 


\markdownInput{0100-enterprise-iot.md}
%\markdownInput{0800-playbooks.md}
%\appendix
%\markdownInput{0900-glossary.md}
%\bibliographystyle{plainnat}
%\bibliography{references}

\end{document}
